\section{Background and Motivation}
\label{sec:background}

\subsection{Distributed Traces}

A distributed trace records the end-to-end execution of a single user
request as it propagates through multiple interacting microservices
\cite{sigelman2010dapper}. Each trace consists of a collection of
\emph{spans}, where every span represents an individual service
invocation together with execution metadata, including the service name,
operation name, start time, duration, and parent-child relationship. The
parent-child relationships among spans allow a trace to be interpreted as
a causal execution structure, from which request paths across services
can be reconstructed. Because traces vary in both execution path and
structural complexity, they provide a rich source of information for
performance analysis and anomaly detection.

\subsection{Service Trace Vectors}

Although distributed traces capture detailed execution behavior, their
variable length and hierarchical structure make them unsuitable for
direct use by conventional machine learning algorithms that require
fixed-dimensional inputs. TraceAnomaly~\cite{liu2020traceanomaly}
addresses this limitation through the Service Trace Vector (STV), which
maps each distributed trace into a fixed-dimensional feature vector while
preserving the correspondence between features and service call paths.

Following TraceAnomaly, each feature 
$j=(\mathrm{service},\mathrm{service\_path})$ corresponds to a unique service reached through a particular root-to-node execution path observed during training.
For a trace \(i\), the feature value is defined as

\begin{equation}
x_{i,j}=\log\!\left(1+d_{i,j}\right),
\label{eq:stv-value}
\end{equation}

where \(d_{i,j}\) denotes the maximum observed duration for feature \(j\)
within trace \(i\). When the same service-path feature appears multiple
times in a trace, the maximum duration is retained. Features that are
absent from a trace are assigned a sentinel value of \(-1\), indicating
that the corresponding service path was not executed. The logarithmic
transformation reduces the dynamic range of span durations while
preserving their relative ordering. Because every feature corresponds to
a specific service path, the resulting representation remains directly
interpretable and enables anomaly scores to be associated with concrete
execution paths rather than latent embedding dimensions.

\subsection{Motivation: Trace-Complexity Bias}
\label{sec:complexity-bias-background}

The STV representation provides an interpretable numerical description
of distributed traces, but it also introduces an important
characteristic: structurally larger traces naturally populate more STV
features than smaller traces. Consequently, traces containing more
services or spans may appear more unusual to an anomaly scoring method,
even when the underlying system behavior is operationally normal.

To investigate whether this structural difference influences anomaly
scoring, we analyze the relationship between baseline STV anomaly scores
and two measures of trace complexity: the number of distinct services
and the total number of spans contained in each trace. The analysis is
performed on real, non-injected test traces collected from three
TrainTicket configurations~\cite{zhou2021trainticket}.

\begin{table}[t]
\centering
\caption{Pearson correlation between baseline STV~+~Isolation Forest anomaly
score and trace-complexity metrics (real, non-injected traces).}
\label{tab:complexity-bias-baseline}
\begin{tabular}{lccc}
\toprule
& \textbf{Dataset 1} & \textbf{Dataset 2} & \textbf{Dataset 3} \\
\midrule
num\_services      & 0.977 & 0.982 & 0.968 \\
num\_spans         & 0.854 & 0.981 & 0.975 \\
total\_duration    & 0.605 & 0.698 & 0.672 \\
max\_span\_duration & 0.550 & 0.627 & 0.681 \\
mean\_span\_duration & 0.035 & $-$0.034 & 0.089 \\
\bottomrule
\end{tabular}
\end{table}


As shown in Table~\ref{tab:complexity-bias-baseline}, baseline anomaly
scores exhibit consistently strong positive correlations with both
complexity measures across all three datasets. Because these
measurements are obtained from traces without injected anomalies, the
observed correlations cannot be attributed to successful detection of
synthetic faults. Instead, they indicate that structural properties of a
trace can substantially influence anomaly scores even when the observed
executions are operationally normal.

Figure~\ref{fig:bias} visualizes this relationship by plotting baseline
anomaly scores against trace size for real, non-injected traces. The
overall trend demonstrates that larger traces tend to receive higher
baseline anomaly scores, motivating the need for an anomaly scoring
strategy that is less sensitive to structural complexity.

\begin{figure}[t]
\centering
% \includegraphics[width=0.45\textwidth]{figures/complexity_bias_scatter.pdf}
\caption{Relationship between baseline STV anomaly score and trace
complexity on real, non-injected traces.}
\label{fig:bias}
\end{figure}

\subsection{Design Intuition}

The observations above suggest that anomaly scoring should depend on how
unusual each feature is relative to its own expected behavior rather
than on its absolute magnitude alone. Features whose latency values are
consistent with their historical behavior should contribute little to the
final anomaly score, even when they occur in a large trace. In contrast,
features exhibiting unusually large deviations from their own learned
normal behavior should receive greater emphasis. This intuition motivates
the Weighted Residual Service Trace Vector (WR-STV) method introduced in
the next section, which preserves the interpretability of the original
STV representation while reducing the influence of trace structural
complexity on anomaly scoring.